\documentclass{article}
\clearpage

\usepackage[utf8]{inputenc}
\usepackage{graphicx} % Required for inserting images
\usepackage{listings}
\usepackage{xcolor}
\usepackage[spanish]{babel}
\usepackage{float}
\usepackage{hyperref}

\title{Proyecto Final GitHub}
\author{Frida Hernández Rodríguez}
\date{Mayo 2025}

\begin{document}

\maketitle
\pagebreak
\tableofcontents
\pagebreak
\section{Introducción}
El presente proyecto tiene como finalidad reforzar los conocimientos y habilidades prácticas en el uso de Git y GitHub como herramientas fundamentales para el control de versiones y la colaboración en entornos de desarrollo. A través de la creación de un repositorio personal, se llevó a cabo la gestión de múltiples ramas y archivos mediante el uso de comandos de Git, fomentando una estructura ordenada y coherente en el manejo de la información.

Cada archivo del proyecto fue desarrollado en una rama distinta y contiene contenido estructurado con buen diseño, incluyendo títulos, subtítulos, listas, tablas, imágenes e hipervínculos. Se realizaron múltiples commits para registrar el proceso de construcción de cada archivo, manteniendo un historial claro de cambios. Posteriormente, se integraron los contenidos a la rama principal (main), resolviendo conflictos de fusión cuando fue necesario.

Asimismo, como parte de la dinámica colaborativa, se utilizaron forks de repositorios de compañeros para practicar el flujo de trabajo de GitHub mediante pull requests. Estas acciones permitieron simular un entorno de colaboración como el que se vive en proyectos reales de desarrollo de software.

Este ejercicio no solo fortalece la comprensión técnica de Git, sino que también desarrolla la capacidad de organización, resolución de problemas y documentación profesional de proyectos.

\href{https://github.com/Fridaaah/ProyectoFinal_GitHub_Frida_Hdz.git}{Repositorio del proyecto en GitHub}


\pagebreak
  \section{Desarrollo}
  
  \subsection{Crear directorio y archivos markdown }
Empezamos creando el directorio del proyecto, despues inicializamos el repositorio y creamos los cuatro archivos markdown (Informacion.md, Proyectos.md, Clasificacion.md y README.md).

\subsubsection{Captura de pantalla}

\begin{figure}[!ht]
\centering
\includegraphics[width=0.7\textwidth]{Crear archivos .md.png}
\caption{Creacion de directorio, archivos y git init}
\end{figure}


\subsection{Primer commit}
Realizamos nuestro primer commit "Crear archivos markdown"

\subsubsection{Captura de pantalla}

\begin{figure}[!ht]
\centering
\includegraphics[width=0.7\textwidth]{Primmer commit.png}
\caption{Primer commit - Crear archivos markdown}
\end{figure}

\pagebreak
\subsection{Creacion de ramas y enlazar repositorio}
Aquí creamos nuestras tres ramas a dicionales a la rama main (informacion, proyectos y clasificacion) y posteriormente enlazamos con nuestro repositorio de GitHub y utilizamos el comando git push con cada una de las ramas creadas.

\subsubsection{Capturas de pantalla}

\begin{figure}[!ht]
\centering
\includegraphics[width=0.7\textwidth]{Creacion de ramas.png}
\caption{Creación de ramas}
\end{figure}

\begin{figure}[!ht]
\centering
\includegraphics[width=0.7\textwidth]{enlazar con repositorio y push de ramas.png}
\caption{Enlace a repositorio y push de ramas}
\end{figure}

\subsection{Trabajar en la rama informacion}
Comenzamos cambiandonos a la rama informacion y empezamos a egregar contenido al archivo informacion.md. Realizamos un total de ocho commits y al final hacemos push a la rama informacion. 

\subsubsection{Capturas de pantalla}

\begin{figure}[!ht]
\centering
\includegraphics[width=0.7\textwidth]{cambiar a rama info.png}
\caption{Cambiar a la rama informacion}
\end{figure}
\pagebreak
\begin{figure}[!ht]
\centering
\includegraphics[width=0.7\textwidth]{1er cambio informacion.png}
\caption{Ejemplo de agregar contenido en informacion.md - rama informacion}
\end{figure}
\begin{figure}[!ht]
\centering
\includegraphics[width=0.7\textwidth]{1er commit info.png}
\caption{Ejemplo de commit}
\end{figure}
\begin{figure}[!ht]
\centering
\includegraphics[width=0.7\textwidth]{git push informacion.png}
\caption{Push a la rama informacion}
\end{figure}
\begin{figure}[!ht]
\centering
\includegraphics[width=0.7\textwidth]{reflejo en repo informacion.png}
\caption{Reflejo del archivo en el repositorio}
\end{figure}
\pagebreak


\subsection{Trabajar en la rama clasificacion}
Comenzamos cambiandonos a la rama clasificacion y empezamos a egregar contenido al archivo Clasificacion.md. sobre el top 6 de mis peliculas favoritas, agregando en formato tabla el título de la película, una imagen, el director y la razón por la cual es de mis peliculas favoritas. Realizamos un total de nueve commits y al final hacemos push a la rama clasificacion. 

\subsubsection{Capturas de pantalla}

\begin{figure}[!ht]
\centering
\includegraphics[width=0.7\textwidth]{cambiar de rama y 1er commit clasificacion.png}
\caption{Cambiar a la rama clasificacion y hacemos el primer commit}
\end{figure}
\begin{figure}[!ht]
\centering
\includegraphics[width=0.7\textwidth]{1er cambio clasificacion.png}
\caption{Ejemplo de agregar contenido en Clasificacion.md - rama clasificacion}
\end{figure}
\begin{figure}[!ht]
\centering
\includegraphics[width=0.7\textwidth]{2do commit clasificacion.png}
\caption{Ejemplo de commit}
\end{figure}

\begin{figure}[!ht]
\centering
\includegraphics[width=0.7\textwidth]{git push clasificacion.png}
\caption{Push a la rama clasificacion}
\end{figure}

\begin{figure}[!ht]
\centering
\includegraphics[width=0.7\textwidth]{reflejo en repo clasificaccion.png}
\caption{Reflejo del archivo en el repositorio}
\end{figure}
\pagebreak


\subsection{Trabajar en la rama proyectos}
Comenzamos cambiandonos a la rama proyectos y empezamos a egregar contenido al archivo Proyectos.md. sobre los proyectos en los que he participado a lo largo de mi trayectoria como estudiante y profesionista. Realizamos un total de seis commits y al final hacemos push a la rama proyectos. 

\subsubsection{Capturas de pantalla}

\begin{figure}[!ht]
\centering
\includegraphics[width=0.7\textwidth]{cambiar de rama y 1er commit proyectos.png}
\caption{Cambiar a la rama proyectos y hacemos el primer commit}
\end{figure}
\begin{figure}[!ht]
\centering
\includegraphics[width=0.7\textwidth]{1er cmabio proyectos.png}
\caption{Ejemplo de agregar contenido en Proyectos.md - rama proyectos}
\end{figure}
\begin{figure}[!ht]
\centering
\includegraphics[width=0.7\textwidth]{2do commit proyectos.png}
\caption{Ejemplo de commit}
\end{figure}

\begin{figure}[!ht]
\centering
\includegraphics[width=0.7\textwidth]{git push proyectos.png}
\caption{Push a la rama proyectos}
\end{figure}

\begin{figure}[!ht]
\centering
\includegraphics[width=0.7\textwidth]{reflejo proyectos.png}
\caption{Reflejo del archivo en el repositorio}
\end{figure}
\pagebreak


\subsection{Trabajar en la rama MAIN}
Comenzamos cambiandonos a la rama MAIN, después traemos el contenido del archivo Proyectos.md localizado en la rama proyectos. Posteriormente realizamos las modificaciones peritinentes y hacemos un total de dos commits. Al final hacemos push a la rama MAIN.
Repetimos los pasos trayendo el contenido del archivo Clasificacion.md de la rama clasifucacion.

\subsubsection{Capturas de pantalla}

\begin{figure}[!ht]
\centering
\includegraphics[width=0.7\textwidth]{cambiar a rama main.png}
\caption{Cambiar a la rama MAIN}
\end{figure}
\begin{figure}[!ht]
\centering
\includegraphics[width=0.7\textwidth]{traer info de proyectos y 1er commit main proyectos.png}
\caption{Traer contenido de Proyecto.md y ejemplo de commit}
\end{figure}
\begin{figure}[!ht]
\centering
\includegraphics[width=0.7\textwidth]{git push main proyectos.png}
\caption{Push a la rama MAIN}
\end{figure}

\begin{figure}[!ht]
\centering
\includegraphics[width=0.7\textwidth]{1er commit main clasificacion.png}
\caption{Traer contenido de Clasificacion.md y ejemplo de commit}
\end{figure}
\begin{figure}[!ht]
\centering
\includegraphics[width=0.7\textwidth]{git push main clasificacion.png}
\caption{Push a la rama MAIN}
\end{figure}


\subsection{Unir las ramas con la rama MAIN}
Unimos la rama informacion con la rama MAIN (automatic merge). Despues unimos la rama proyectos con la rama MAIN y solucionamos el conflicto que nos aparece de manera manual, al final hacemos el commit y el rerspetivo push a la rama MAIN. Este proceso se repite con la rama clasificacion.

\subsubsection{Capturas de pantalla}

\begin{figure}[!ht]
\centering
\includegraphics[width=0.7\textwidth]{merge informacion con main.png}
\caption{Unimos la rama informacion con la rama MAIN}
\end{figure}
\begin{figure}[!ht]
\centering
\includegraphics[width=0.7\textwidth]{merge proyectos con main CONFLICT.png}
\caption{Unimos rama proyectos con MAIN pero hay CONFLICT}
\end{figure}
\begin{figure}[!ht]
\centering
\includegraphics[width=0.7\textwidth]{CONFLICT proyectos y main.png}
\caption{Ejemplo del CONFLICT}
\end{figure}
\begin{figure}[!ht]
\centering
\includegraphics[width=0.7\textwidth]{Accept current changes proyectos.png}
\caption{Ejemplo de resolver el CONFLICT aceptando los cambios actuales}
\end{figure}
\begin{figure}[!ht]
\centering
\includegraphics[width=0.7\textwidth]{git commit el merge de proyectos a main.png}
\caption{Ejemplo de commit de la union de las dos ramas}
\end{figure}
\begin{figure}[!ht]
\centering
\includegraphics[width=0.7\textwidth]{CONFLICT proyectos y main.png}
\caption{Ejemplo de push para visualizar las ramas unidas}
\end{figure}

\subsection{Hacer 5 forks y pull requests}
Forkeamos 5 repositorios de nuestros compañeros y a cada uno le realizamos un pull request con la modificacion que realizamos en su archivo Clasificacion.md. 

\subsubsection{Capturas de pantalla}

\begin{figure}[!ht]
\centering
\includegraphics[width=0.7\textwidth]{forks y pull req.png}
\caption{Forks y Pull Requests}
\end{figure}
\pagebreak


\section*{Conclusión}

Este proyecto permitió aplicar de forma práctica los conocimientos sobre el manejo de control de versiones con \texttt{Git} y la plataforma \texttt{GitHub}, mediante la creación de un repositorio con múltiples ramas, commits estructurados y fusión de contenidos. Además, se fortalecieron habilidades como el trabajo ordenado, la documentación de procesos y la colaboración a través de pull requests. La experiencia brindó una comprensión más profunda de las herramientas de desarrollo modernas utilizadas en entornos profesionales.


\end{document}
